\section{Database Design}
\label{sec:sd-database}

The schema has 24 tables in five groups (Table~\ref{tab:sd-tables}). Figure~\ref{fig:sd-domain-model} shows the main persistent entities and their associations as a UML class diagram; tables that only support these entities (sessions, settings, revisions and exports) are left out.

\begin{figure}[H]
    \centering
    \resizebox{\linewidth}{!}{% UML class diagram of the main persistent entities (Section 5.7).
\begin{tikzpicture}[uml]
    \begin{class}[text width=2.9cm]{User}{0,0}
        \attribute{id : UUID}
        \attribute{email : String}
        \attribute{display\_name : String}
    \end{class}
    \begin{class}[text width=3.1cm]{Project}{5.2,0}
        \attribute{id : UUID}
        \attribute{name : String}
        \attribute{initial\_message\_id : UUID [0..1]}
    \end{class}
    \begin{class}[text width=2.9cm]{SourceDocument}{10.4,0}
        \attribute{id : UUID}
        \attribute{title : String}
        \attribute{text : String}
    \end{class}
    \begin{class}[text width=2.9cm]{Thread}{0,-3.4}
        \attribute{id : UUID}
        \attribute{title : String}
        \attribute{current\_stage : String}
    \end{class}
    \begin{class}[text width=2.9cm]{Message}{0,-7.6}
        \attribute{id : UUID}
        \attribute{role : MessageRole}
        \attribute{content : String}
        \attribute{delivery : DeliveryStatus}
    \end{class}
    \begin{class}[text width=3.1cm]{Run}{5.2,-3.4}
        \attribute{id : UUID}
        \attribute{status : RunStatus}
        \attribute{workflow\_name : String}
        \attribute{workflow\_version : Integer}
        \attribute{run\_config : JSON}
    \end{class}
    \begin{class}[text width=3.1cm]{RunStep}{3.9,-7.6}
        \attribute{sequence : Integer}
        \attribute{node : String}
        \attribute{agent\_role : String [0..1]}
        \attribute{usage : Usage}
    \end{class}
    \begin{class}[text width=2.9cm]{RunEvent}{7.3,-7.6}
        \attribute{sequence : Integer}
        \attribute{type : String}
        \attribute{payload : JSON}
    \end{class}
    \begin{class}[text width=2.9cm]{Interrupt}{10.7,-7.6}
        \attribute{node : String}
        \attribute{options : JSON}
        \attribute{status : InterruptStatus}
    \end{class}
    \begin{class}[text width=2.9cm]{HumanDecision}{10.7,-10.8}
        \attribute{option : String}
        \attribute{feedback : String [0..1]}
        \attribute{answers : JSON}
    \end{class}
    \begin{class}[text width=3.0cm]{Artifact}{10.4,-3.4}
        \attribute{kind : String}
        \attribute{version : Integer}
        \attribute{review\_status : ReviewStatus}
        \attribute{content : JSON}
    \end{class}
    \begin{class}[text width=3.0cm]{ArtifactItem}{14.6,-3.4}
        \attribute{kind : String}
        \attribute{key : String}
        \attribute{version : Integer}
        \attribute{content : JSON}
    \end{class}
    \begin{class}[text width=3.0cm]{TraceLink}{14.6,-7.6}
        \attribute{link\_type : String}
        \attribute{source\_quote : String [0..1]}
        \attribute{target\_quote : String [0..1]}
    \end{class}

    \association{User}{members}{*}{Project}{projects}{*}
    \composition{Project}{}{*}{SourceDocument}
    \composition{Project}{}{*}{Thread}
    \composition{Thread}{}{*}{Message}
    \composition{Thread}{}{*}{Run}
    \composition{Run}{}{*}{RunStep}
    \composition{Run}{}{*}{Interrupt}
    \association{Interrupt}{}{1}{HumanDecision}{}{0..1}
    \composition{Project}{}{*}{Artifact}
    \composition{Artifact}{}{*}{ArtifactItem}
    \composition{Run}{}{*}{RunEvent}
    \draw[thick] (ArtifactItem.north) -- ++(0,0.5) -| node[pos=0.25, above, font=\tiny] {parent 0..1} ($(ArtifactItem.north)+(-0.9,0)$);
    \draw[thick, ->] ($(TraceLink.north)+(-0.5,0)$) -- node[left, font=\tiny] {source 1} ($(ArtifactItem.south)+(-0.5,0)$);
    \draw[thick, ->] ($(TraceLink.north)+(0.5,0)$) -- node[right, font=\tiny] {target 1} ($(ArtifactItem.south)+(0.5,0)$);
    \node[font=\tiny] at ($(TraceLink.north)+(-0.75,0.25)$) {*};
    \node[font=\tiny] at ($(TraceLink.north)+(0.75,0.25)$) {*};
\end{tikzpicture}
}
    \caption{UML class diagram of the main persistent entities.}
    \label{fig:sd-domain-model}
\end{figure}

\begingroup\setstretch{1.15}
\begin{longtable}{|p{0.2\linewidth}|p{0.68\linewidth}|}
    \caption{Groups of tables of the database.}
    \label{tab:sd-tables} \\
    \hline
    \textbf{Group} & \textbf{Tables and what they hold} \\
    \hline
    \endfirsthead
    \hline
    \textbf{Group} & \textbf{Tables and what they hold} \\
    \hline
    \endhead
    Identity and settings (4) & Users, with their monthly budget; their login sessions; their settings, such as the theme, the default output language and the default run settings; and the model chosen for each agent role. \\ \hline
    Workspace (4) & Projects with their owner; the counters that number the keys of the items of a project; the uploaded source documents and their versions. \\ \hline
    Conversation and execution (9) & Threads and their messages, with the delivery status of each message; runs with their status, workflow version and frozen run configuration; the steps of each run with their model, usage and saved state; the ordered event log of each run; the review requests and the human decisions that answer them; the directives that the intent router and the rework control write for the next node; and the usage of the model calls that classify a message before any run exists. \\ \hline
    Artifacts and trace (5) & Artifacts and their revisions; artifact items and their revisions; and the typed trace links between items, with their quotes. \\ \hline
    Supporting (2) & Exports with the versions of the artifacts they hold; and the keys that make repeated requests idempotent. \\ \hline
\end{longtable}
\endgroup

The design associates users with projects as members (Figure~\ref{fig:sd-domain-model}); the current implementation gives each project a single owner, which is why FR-03 is partly met.

\subsection{Artifact Storage}
\label{sec:sd-artifact-storage}
Requirements engineering output has two levels, each in its own generic table: a table of artifacts and a table of artifact items, with the fields of each kind stored in a JSON content field validated against a versioned content schema. The items and the trace links between them form a property graph.

The main design decisions and their trade-offs are summarized in Table~\ref{tab:sd-db-decisions}.

\begin{table}[H]
    \centering
    \caption{Main decisions of the database design.}
    \label{tab:sd-db-decisions}
    \begin{tabularx}{\linewidth}{|p{2.6cm}|X|X|X|}
        \hline
        \textbf{Question} & \textbf{Options considered} & \textbf{Decision} & \textbf{Trade-off} \\
        \hline
        Fields of the artifact and item kinds & One table per kind with typed columns; generic tables with a JSON content column & Generic tables, one per level, with a JSON content field validated against a versioned content schema & A field change is a schema change in code, not a migration, which suits fields that change while prompts are tuned; the database does not check the fields, so the application must \\
        \hline
        Names of the workflow (roles, nodes, options, kinds, link types) & Database enums; lookup tables; text checked against the workflow definition & Text checked against the versioned workflow definition; every run stores its workflow name and version & A new experiment is a new definition version with no migration, and old runs stay readable; every write must be checked by the application \\
        \hline
        Storage of the trace graph & Graph database; relational database & Relational database & The rest of the data (users, runs, events, checkpoints) is relational and the agent runtime can store its checkpoints there; the graph queries needed (trace of an item, matrix, coverage) are shallow \\
        \hline
    \end{tabularx}
\end{table}

The design targets PostgreSQL with the pgvector extension for the embeddings of items. The current implementation runs on SQLite through the same SQLAlchemy models and Alembic migrations, which keeps the development setup to one file; Section~\ref{sec:impl-backend} describes it.
